% ---------------------------------------------------------------------------
% Tables from docs/methods_draft.md, in LaTeX.
%
% Preamble (add once to your document):
%
%   \usepackage{booktabs}
%   \usepackage{tabularx}
%   \newcolumntype{L}[1]{>{\raggedright\arraybackslash\hsize=#1\hsize}X}
%
% The L{w} widths are relative weights that must SUM TO the number of X
% columns in that table (4 for a four-column table, and so on). Shift weight
% between them if a column looks cramped; keep the total unchanged.
%
% Two-column (IEEE-style) documents: the wide tables here -- 1, 3 and 6 --
% will be tight in one column. Use table* and \textwidth to span both.
% ---------------------------------------------------------------------------


% === 1. Dataset selection against requirements (Section 2) =================
\begin{table}[t]
  \centering
  \caption{Candidate datasets against the requirements the system imposes.
           Only CroHD satisfies all five.}
  \label{tab:dataset-requirements}
  \begin{tabularx}{\linewidth}{@{}L{0.80}L{1.40}L{0.95}L{0.85}@{}}
    \toprule
    \textbf{Requirement} & \textbf{Why it is necessary}
      & \textbf{SCUT-HEAD} & \textbf{CroHD} \\
    \midrule
    Heads, not bodies
      & bodies are occluded in dense crowds
      & yes & yes \\
    Camera above the crowd
      & determines whether faces or backs of heads are seen
      & partly --- elevated but front-facing
      & \textbf{yes} \\
    Genuinely dense scenes
      & the regime of interest
      & no --- tens per frame
      & \textbf{yes --- hundreds} \\
    Box annotations
      & point labels cannot train a detector
      & yes & \textbf{yes} \\
    Video with identities
      & required for velocity and crowd pressure
      & no --- stills
      & \textbf{yes} \\
    \bottomrule
  \end{tabularx}
\end{table}


% === 2. Delivered model configuration (Section 3) ==========================
\begin{table}[t]
  \centering
  \caption{Configuration of the delivered detector.}
  \label{tab:delivered-model}
  \begin{tabularx}{\linewidth}{@{}lX@{}}
    \toprule
    Initialisation & COCO-pretrained YOLOv8n \\
    Training data  & CroHD, 1{,}187 images, 253{,}376 head boxes \\
    Classes        & 1 (\texttt{head}) \\
    Resolution     & 1280 \\
    Epochs         & 10, batch 2 \\
    Augmentation   & horizontal flip, $\pm$30\% scale; mosaic, rotation, shear
                     and perspective disabled \\
    Parameters     & 3{,}138{,}003 \\
    Inference      & 30.8\,ms/frame at \texttt{imgsz=1280} (RTX 3050 6\,GB),
                     against 143.4\,ms on CPU \\
    \bottomrule
  \end{tabularx}
\end{table}


% === 3. The train / validation / test split (Section 4) ====================
\begin{table}[t]
  \centering
  \caption{The three data roles. Only ShanghaiTech played no part in either
           fitting or model selection, so only it is a genuine test set.}
  \label{tab:split}
  \begin{tabularx}{\linewidth}{@{}L{0.60}L{0.95}L{1.20}L{1.25}@{}}
    \toprule
    \textbf{Role} & \textbf{Data} & \textbf{Size} & \textbf{Purpose} \\
    \midrule
    Train
      & HT21-01, 02, 03
      & 1{,}187 images / 253{,}376 boxes
      & fits the weights \\
    Validation
      & HT21-04 (whole scene)
      & 100 images / 17{,}665 boxes
      & \textbf{selects the retained checkpoint} \\
    \textbf{Test}
      & \textbf{ShanghaiTech Part A}
      & \textbf{482 images / 241{,}275 points}
      & \textbf{independent evaluation} \\
    \bottomrule
  \end{tabularx}
\end{table}


% === 4. Cross-domain count ratio (Section 4) ===============================
\begin{table}[t]
  \centering
  \caption{Count ratio on 482 ShanghaiTech images, by crowd density. All three
           models scored identically at confidence 0.25, \texttt{imgsz=1024},
           detection cap 4{,}000.}
  \label{tab:count-ratio}
  \begin{tabular}{@{}lrrr@{}}
    \toprule
    \textbf{heads/image} & \textbf{SCUT-HEAD A} & \textbf{SCUT-HEAD A+B}
      & \textbf{CroHD (delivered)} \\
    \midrule
    100--200  & 0.269 & 0.421 & \textbf{0.975} \\
    200--400  & 0.180 & 0.288 & \textbf{0.740} \\
    1200+     & 0.012 & 0.024 & \textbf{0.204} \\
    \midrule
    usable ceiling (ratio $\geq$ 0.8) & never & never
      & \textbf{200 heads/image} \\
    \bottomrule
  \end{tabular}
\end{table}


% === 5. Standard detection metrics (Section 4) =============================
\begin{table}[t]
  \centering
  \caption{Detection metrics on the held-out CroHD scene at the deployed
           confidence of 0.25. HT21-04 is the delivered model's validation
           scene, so its row is an upper bound; for the two controls these
           images are unseen.}
  \label{tab:detection-metrics}
  \begin{tabular}{@{}lrrrrr@{}}
    \toprule
    \textbf{Model} & \textbf{Precision} & \textbf{Recall} & \textbf{F1}
      & \textbf{mAP50} & \textbf{mAP50-95} \\
    \midrule
    SCUT-HEAD A    & 0.314 & 0.065 & 0.107 & 0.028 & 0.006 \\
    SCUT-HEAD A+B  & 0.535 & 0.328 & 0.406 & 0.213 & 0.041 \\
    \textbf{CroHD (delivered)}
      & \textbf{0.747} & \textbf{0.696} & \textbf{0.720}
      & \textbf{0.631} & \textbf{0.202} \\
    \bottomrule
  \end{tabular}
\end{table}


% === 6. The two decision rules (Section 5) =================================
\begin{table}[t]
  \centering
  \caption{The two decision rules. They use different level vocabularies
           because they make different claims: one is absolute, the other is
           relative to the scene it is watching.}
  \label{tab:decision-rules}
  \begin{tabularx}{\linewidth}{@{}L{0.65}L{1.20}L{1.15}@{}}
    \toprule
      & \textbf{calibrated rule} & \textbf{scene-relative rule} \\
    \midrule
    Implemented in
      & \texttt{src/cdw/warning.py}
      & \texttt{scripts/demo\_uncalibrated.py} \\
    Levels
      & NORMAL / WATCH / \textbf{WARNING}
      & LEARNING / WATCH / \textbf{ELEVATED} \\
    Signals
      & density in a crowded Fruin band \textbf{and} crowd pressure rising
      & occupancy above this scene's own normal \textbf{and} net inflow
        positive \\
    Units
      & people/m\textsuperscript{2} --- absolute
      & people --- relative to this camera \\
    Thresholds from
      & cited literature (Fruin) plus a swept parameter
      & the footage itself \\
    Claims
      & conditions resemble values documented before past flow breakdowns
      & conditions are unusual \textbf{for this camera} \\
    \bottomrule
  \end{tabularx}
\end{table}
